\chapter{MATERIAIS E MÉTODOS}
\label{capitulo3}

Descreva neste capítulo os materiais, ferramentas, bases de dados e procedimentos metodológicos utilizados para conduzir a pesquisa, garantindo rigor científico e reprodutibilidade dos experimentos ou coletas de dados.

\section{MATERIAIS E BASES DE DADOS}
\label{sec:materiais}

Descreva os equipamentos, softwares, bases de dados secundárias e insumos empregados. Indique detalhadamente fontes, períodos temporais de coleta e características técnicas das ferramentas.

\section{MÉTODOS E PROCEDIMENTOS}
\label{sec:metodos}

Apresente as etapas do método de pesquisa, as técnicas de amostragem, modelagem matemática ou simulação adotadas.

\subsection{Padronização de Tabelas no PPGT}
\label{subsec:tabelas}

Diferente dos \textit{Quadros} (que possuem bordas laterais fechadas e abrigam dados predominantemente textuais), as \textbf{Tabelas} são destinadas a dados numéricos e seguem a norma de apresentação tabular do IBGE e do PPGT:
\begin{itemize}
  \item Bordas laterais abertas (sem linhas verticais nas extremidades esquerda e direita);
  \item Linhas horizontais no topo, abaixo dos cabeçalhos das colunas e na base inferior da tabela;
  \item Cabeçalhos de colunas em negrito;
  \item Indicação da fonte na base inferior com o comando \texttt{\textbackslash fonte\{...\}} quando extraída ou adaptada de outros trabalhos (dispensável quando elaborada pelo próprio autor).
\end{itemize}

A Tabela~\ref{tab:exemplo_tabela} ilustra a formatação correta de uma tabela segundo o padrão exigido:

\begin{table}[htbp]
  \centering
  \caption{Exemplo de tabela numérica formatada segundo as normas do PPGT}
  \label{tab:exemplo_tabela}
  \small
  \begin{tabular}{lrr}
    \hline
    \textbf{Categoria de Análise} & \textbf{Frequência Absoluta} & \textbf{Percentual (\%)} \\
    \hline
    Classe A & 142 & 45,8 \\
    Classe B & 98 & 31,6 \\
    Classe C & 70 & 22,6 \\
    \hline
    \textbf{Total} & \textbf{310} & \textbf{100,0} \\
    \hline
  \end{tabular}
  % Se a tabela for elaborada integralmente pelo próprio autor, remova ou comente a linha \fonte abaixo:
  \fonte{adaptado de \citet{dnit2006}}
\end{table}

\subsection{Regra Específica do PPGT para o 3º Nível de Subitens}
\label{subsec:terceiro_nivel}

Uma diretriz exclusiva das normas do PPGT é o tratamento do 3º nível de subseção (\texttt{\textbackslash subsubsection}):
\begin{itemize}
  \item No texto, o título do 3º nível aparece em formato itálico com identificador alfabético: \textbf{(a)}, \textbf{(b)}, \textbf{(c)}... em vez da numeração decimal clássica (ex.: 3.2.1);
  \item No Sumário (Índice), os títulos de 3º nível não são exibidos (a profundidade é limitada a 2 níveis, conforme exigido institucionalmente).
\end{itemize}

\subsubsection{Primeira Etapa do Procedimento Analítico}

Exemplo de texto pertencente ao primeiro subitem do 3º nível, identificado automaticamente pelo prefixo (a).

\subsubsection{Segunda Etapa do Procedimento Analítico}

Exemplo de texto pertencente ao segundo subitem do 3º nível, identificado automaticamente pelo prefixo (b).
